\documentclass{article}
\usepackage{hyperref}
\usepackage{Style}

\nocite{*} % Comentar si quiero citar
%\addbibresource{bibliografia.bib} % Quitar el comentado si quiero usar bibliografia

\begin{document}

\begin{minipage}{2.5cm}
    \includegraphics[width=2cm]{imagen_puc.jpg}
\end{minipage}
\begin{minipage}{12cm}
    {\sc Pontificia Universidad Católica de Chile\\
    Facultad de Matemáticas\\
    Departamento de Matemática\\
    Profesor: Manuel Cabezas -- Ayudante: Benjamín Mateluna}
\end{minipage}
\vspace{2ex}

{\centerline{\bf Introducción a la Geometría - MAT1304}
\centerline{\bf Ayudantía 25 - Repaso I7}}
\centerline{\bf 18 de junio de 2026}

\vspace{1cm}
\noindent\textbf{Problema 1.} Sean $f,g\in k[x,y]$, demuestre que
\begin{enumerate}
    \item El conjunto de ceros de $fg$ es igual a la unión del conjunto de ceros de $f$ y $g$.
    \item Si $k=\R$, el conjunto de ceros de $f^{2}+g^{2}$ es igual a la intersección del conjunto 
    de ceros de $f$ y $g$.
\end{enumerate}

\vspace{5mm}
\noindent\textbf{Problema 2.} Dada la circunferencia $x^{2}+y^{2}-6x-2y+6=0$, determine los 
valores de $m$ para los cuales la recta $y=mx+3$:
\begin{enumerate}
    \item Intersecte a la circunferencia en dos puntos.
    \item Sea tangente a la circunferencia.
    \item No tenga puntos en común con la circunferencia.
\end{enumerate}

\vspace{5mm}
\noindent\textbf{Problema 3.} Encuentre la cónica que pasa por los puntos 
$(0,1), (2,1), (-1,0), (0,2)$ y $(1,0)$.

\vspace{5mm}
\noindent\textbf{Problema 4.} Describa el conjunto solución de la ecuación 
$5x^{2}+4xy+2y^{2}-24x-12y+29=0$.

\vspace{1cm}
\noindent\textbf{\underline{Ejercicios Propuestos:}}

\vspace{5mm}
\noindent\textbf{Extra 1.} Encontrar las rectas tangentes a $x^{2}+y^{2}=1$ que son paralelas a la
recta de ecuación $y=-2x+5$.

\vspace{5mm}
\noindent\textbf{Extra 2.} Hallar e identificar la ecuación del lugar geométrico de los puntos
tales que su distancia a la recta $y=-8$ es igual al doble de su distancia con el punto $(0,-2)$.

%\printbibliography % Quitar el comentado si quiero usar bibliografia

\end{document}
